\documentclass{article}

% If you need to pass options to natbib, use, e.g.:
%\PassOptionsToPackage{numbers, compress}{natbib}
% before loading staix_2026.

% ============================================================
% STAI-X 2026 track options
% Default: anonymous submission, Full Paper Track
% ============================================================
\usepackage{staix_2026}
% \usepackage[poster]{staix_2026}        % anonymous submission, Poster Track
% \usepackage[final]{staix_2026}         % camera-ready, Full Paper Track
% \usepackage[poster,final]{staix_2026}  % camera-ready, Poster Track
% \usepackage[preprint]{staix_2026}      % non-anonymous preprint
% \usepackage[poster,preprint]{staix_2026}

\usepackage[utf8]{inputenc}
\usepackage[T1]{fontenc}
\usepackage{hyperref}
\usepackage{url}
\usepackage{booktabs}
\usepackage{amsfonts}
\usepackage{nicefrac}
\usepackage{microtype}
\usepackage{xcolor}
\usepackage{graphicx}
\usepackage{amsmath,amssymb,amsthm}

\title{STAI-X 2026 Submission Title}

\author{%
  First Author \\
  Affiliation \\
  Address \\
  \texttt{first.author@example.edu}
  \And
  Second Author \\
  Affiliation \\
  Address \\
  \texttt{second.author@example.edu}
}

\begin{document}

\maketitle

\begin{abstract}
This document provides the author shell for STAI-X 2026 submissions. Replace the title, author block, abstract, and main text with your own content. The active track and page-limit policy are controlled by the style-file options.
\end{abstract}

\staixcontentlimitnotice

\section{Introduction}
This file is the starting point for authors preparing a submission to \@staixshortname. By default, the template compiles an anonymous submission for the Full Paper Track. To prepare a Poster Track submission, switch to
\verb|\usepackage[poster]{staix_2026}| in the preamble.

The template follows the overall design philosophy of recent NeurIPS and ICLR templates while adding explicit support for STAI-X 2026's two-track structure:
\begin{itemize}
    \item \textbf{Full Paper Track:} Up to 8 pages of main content including figures and tables, excluding references.
    \item \textbf{Poster Track:} Up to 2 pages of main content including figures and tables, excluding references.
\end{itemize}

\section{Formatting and style}
Authors should not modify the margins, font sizes, spacing, or title block defined in the style file. The style file controls anonymity, first-page footer text, and line numbering for the review version.

\subsection{Citations}
Use standard citation commands from \texttt{natbib}, such as \verb|\citet| and \verb|\citep|. For example, \citet{lamport1994latex} provides a classical reference on \LaTeX{}, while modern machine-learning conferences often use variants of the same submission-style architecture \citep{goodfellow2016deep}.

\subsection{Figures and tables}
Place figure captions below figures and table captions above tables.

\begin{figure}[h]
\centering
\fbox{\rule[-.5cm]{0cm}{4cm} \rule[-.5cm]{4cm}{0cm}}
\caption{Sample figure caption.}
\label{fig:sample}
\end{figure}

\begin{table}[h]
\caption{Sample table caption.}
\label{tab:sample}
\centering
\begin{tabular}{lc}
\toprule
Method & Accuracy \\
\midrule
Method A & 0.90 \\
Method B & 0.93 \\
\bottomrule
\end{tabular}
\end{table}

\section{Theorem-like environments}
The template does not force theorem environments, but authors can define them in the usual way.

\newtheorem{theorem}{Theorem}
\begin{theorem}
Suppose the style file is used as intended. Then the compiled manuscript conforms to the STAI-X 2026 review format.
\end{theorem}

\section{Acknowledgments}
For anonymous submissions, the \texttt{ack} environment is hidden automatically. It reappears in final and preprint modes.

\begin{ack}
We thank the STAI-X 2026 organizers for creating a venue at the interface of statistics and AI.
\end{ack}

\bibliographystyle{plainnat}
\bibliography{staix_2026_sample}

\appendix
\section{Appendix}
Additional proofs, implementation details, and supplementary experiments may be placed after the references unless the call for papers specifies otherwise.

\end{document}
